\chapter{Structured Reporting in Radiology}

\label{ch:ch36_structured_reporting}

\section{Purpose}
Structured reporting improves completeness, consistency, communication, and data reuse when the structure matches the clinical task. It should support expert judgment rather than force findings into an inappropriate template.

\begin{pearlbox}[Clinical Pearl: The referring clinician's perspective]
Surgeons and physicians often scan the impression first, then look at specific findings. A well-structured impression that leads with the most clinically important finding and ends with a clear recommendation saves time and reduces errors in patient management.
\end{pearlbox}

\section{Essential Elements}
Every radiology report should contain the following elements:

\begin{table}[h]
\centering
\caption{Essential elements of a radiology report.}
\label{tab:report_elements}
\begin{tabular}{p{3.5cm} p{10cm}}
\toprule
\textbf{Element} & \textbf{Content} \\
\midrule
Indication & The clinical question or reason for the examination; brief relevant history \\
Comparison & Prior studies available for comparison, with dates \\
Technique & Modality, protocol, contrast agent and dose, reconstruction parameters \\
Findings & Systematic description of all relevant normal and abnormal observations \\
Limitations & Technical factors that reduce diagnostic confidence \\
Impression & Ranked summary of the most clinically important findings, with calibrated confidence \\
Recommendations & Specific next steps when appropriate \\
\bottomrule
\end{tabular}
\end{table}

\section{Templates by Examination}
Use a consistent template for each examination type. Below are representative structures (abbreviated for space):

\subsection{Chest CT}
Indication, comparison, technique (contrast yes/no, phase). Findings organized by: lungs and airways, pleura, mediastinum and hila, heart and pericardium, upper abdomen, bones and soft tissues. Impression with ranked findings. Recommendation.

\subsection{Abdominal CT}
Indication, comparison, technique, contrast details. Findings organized by organ: liver, biliary system, pancreas, spleen, kidneys and adrenals, GI tract, mesentery, vascular, bones and soft tissues. Impression. Recommendation.

\subsection{MRI Brain}
Indication, comparison, technique (sequences: T1, T2, FLAIR, DWI, SWI, post-gadolinium). Findings organized by: grey--white differentiation, ventricles and sulci, midline structures, brainstem and cerebellum, extra-axial spaces, orbits. Impression. Recommendation.

\section{RADS Lexicons}
Standardised lexicons improve communication and guide management. They should be used only within their intended scope.

\begin{table}[h]
\centering
\caption{Standardised RADS lexicons and their applications.}
\label{tab:radss}
\begin{tabular}{p{2.5cm} p{3.5cm} p{3cm} p{4.5cm}}
\toprule
\textbf{Lexicon} & \textbf{Organ / Modality} & \textbf{Score Range} & \textbf{Key Threshold} \\
\midrule
BI-RADS & Breast (mammography, ultrasound [US], MRI) & 0--6 & BI-RADS 4 requires tissue sampling; BI-RADS 5 highly suggestive \\
Bosniak & Renal cysts (CT/MRI) & I, II, IIF, III, IV & Bosniak III--IV requires surgical or urological consideration \\
LI-RADS & Liver (hepatocellular carcinoma [HCC] surveillance in cirrhosis) & LR-1 to LR-5 & LR-5 = definite HCC; LR-M = non-HCC malignancy \\
Lung-RADS & Lung cancer screening CT & 0, 1, 2, 3, 4A, 4B, 4X & Category 4A/4B/4X triggers diagnostic workup \\
PI-RADS & Prostate (multiparametric MRI [mpMRI]) & 1--5 & PI-RADS 4 and 5 warrant targeted biopsy \\
TI-RADS & Thyroid (US) & TR1--TR5 & Points-based; determines fine-needle aspiration (FNA) threshold \\
\bottomrule
\end{tabular}
\end{table}

\section{Impression and Communication}
Lead with the clinically important conclusion. Distinguish urgent findings, uncertainty, and actionable follow-up. Document communication of critical results according to local policy and protect patient identifiers in educational or research copies.

\subsection{Communicating Critical Findings}
When a critical finding is identified:
\begin{enumerate}
    \item Document the finding in the report with clear, unambiguous language.
    \item Contact the referring clinician directly by telephone or secure messaging.
    \item Record the name of the person informed, the time of communication, and any agreed plan.
    \item Follow up to ensure the result was received and acted upon.
\end{enumerate}

\begin{pearlbox}[Clinical Pearl: Recommendations must be specific]
``Consider follow-up'' is unhelpful. State the modality, the timeframe, and the specific question the follow-up is designed to answer: ``Recommend chest CT in 3 months to assess interval stability of the right upper lobe 8 mm solid nodule (Lung-RADS 3).''
\end{pearlbox}

\begin{pearlbox}[Impression:]
A structured report ensures completeness, clarity, and clinical utility. Use standardised templates, validated RADS lexicons where applicable, calibrated language for uncertainty, and specific actionable recommendations. Always document communication of critical findings.
\end{pearlbox}

\clearpage
